\einheitenkopf{Einheit 1 von 5 · Proportionale Funktion}
\zweigzeile{Hier lernst du, Funktionen mit $y = m \cdot x$ zu berechnen, zu zeichnen und zu erkennen · neu in diesem Jahr · P10 oft · baut auf: Koordinatensystem, proportionale Zuordnungen (Dreisatz)}
\erg{1}{a) $y = 1$, Gerade durch $(0|0)$ und $(1|1)$ \quad b) $y = 3$, durch $(0|0)$ und $(1|3)$ \quad c) $y = 4$, durch $(0|0)$ und $(1|4)$ \quad d) $y = 5$, durch $(0|0)$ und $(1|5)$}
\erg{2}{e) $y = 5$, Gerade durch $(0|0)$ und $(2|5)$ \quad f) z.\,B. $x = 3$, $y = 1$, Gerade durch $(0|0)$ und $(3|1)$ \quad g) $y = -2$, Gerade durch $(0|0)$ und $(1|{-2})$, sie fällt}
\erg{3}{a) $y = 3$ \quad b) $y = -3$ \quad c) $x = 4$ \quad d) $1{,}5 \cdot 2 = 3$, stimmt \quad e) $y = 1{,}5 \cdot 12 = 18$ \quad f) $2 \to 3$, $\cdot 6$: $12 \to 18$, stimmt}
\erg{4}{a) ja, $y = 6x$ \quad b) nein, $14 : 4 = 3{,}5$, aber $20 : 6 \approx 3{,}33$ \quad c) nein, $5 : 1 = 5$, aber $8 : 2 = 4$ \quad d) ja, $y = -1{,}5x$ \quad e) ja, $y = 1{,}8x$ \quad f) nein, bei 0 km kostet es schon 4 €}
\erg{5}{a) Tom hat subtrahiert statt dividiert; der Faktor ist der Quotient $m = 21 : 3 = 7$ (auch $42 : 6 = 7$), $y = 7x$ \quad b) $m = 18 : 4 = 4{,}5$, $y = 4{,}5x$}
\erg{6}{a) Für $x = 0$ ist $y = m \cdot 0 = 0$, also liegt $(0|0)$ immer auf dem Graphen. \quad b) nein: Bei $m = -4$ gehört zu einem positiven $x$ immer ein negatives $y$, $P$ hat aber $y = 5 > 0$.}
\erg{7}{a) $y = 40x$ \quad b) Gerade durch $(0|0)$ und $(10|400)$ \quad c) In jeder Minute kommen 40 Liter dazu; 200 Liter nach 5 Minuten.}
